\BookPart{序}
岁在丙午，时维二〇二六年。

是岁也，九州风起，四海云兴。中国之人，似于一夕之间，忽寻得失落千载之根，洞悉湮没万世之史。昔之所学，尽成疑案；旧之所闻，皆为谬传。山川易位，日月改悬，古今帝王之世系，东西万国之源流，皆有新解。天下之大，竟无一物不可溯其本，无一邦不可考其宗。

于是奇说蜂起，异论纷纭。有人指欧罗巴为故秦之疆，有人谓美利坚乃九州之壤；阿尔卑斯，实为峨眉；多瑙长河，本是楚水。金字塔者，华夏先王之陵寝也；罗马古城者，秦人西迁之遗都也。至若成吉思汗、忽必烈、耶稣、亚里士多德之属，亦皆与大明宗室牵连。千年之人物，万里之山河，纵隔重洋，虽悬百世，苟得一音相近、一字相通，便可认祖归宗，重入华夏版图。

更有博闻之士，穷究造化，旁通鬼神。谓天下之人，原非同源；华夏独承天命，异邦别有根由。乃援《山海经》以证万国，取《红楼梦》以解兴亡，考《西游记》而识诸夷，观西洋油画而辨人兽。凡形貌之异、风俗之殊，皆可推及上古；凡史籍之阙、文物之残，悉谓后人所毁。其说愈奇，其信愈笃；其证愈寡，其辞愈坚。

或问曰：「如此恢宏旧事，何以历代史官俱未尝载？」

答曰：「史已遭篡。」

又问：「何以天下典籍、金石、舆图、碑铭，皆与此说不合？」

答曰：「皆已遭篡。」

再问：「何人能尽篡天下之史？」

则曰：「天下之人，共谋之耳。」

至此，疑者无从置喙，信者益觉豁然。盖史之可信与否，不在考据之精粗，而在觉醒之先后。信之则为明哲，不信则为蒙昧；异议愈众，愈足见其秘密之深；反证愈多，愈足彰其阴谋之巨。此诚千古未有之奇观也。

昔太史公周览天下，访遗老，察故迹，参诸异闻，而后著成一家之言。今之诸贤则不然，端坐斗室，执方寸之屏，纵横古今，驰骋八荒。顷刻之间，可辨万年真伪；谈笑之际，能复四海舆图。既不必穷年累月以读史，亦无须跋山涉水以访古。一图足以易天下之疆，一谐足以改千年之史，一念足以破万国之谜。

呜呼！昔人苦求之真，今人举手可得；昔人不敢言之事，今人无所不言。秦皇汉武未尝知其疆域之广，唐宗宋祖未尝晓其功业之隆，明太祖、成祖泉下有知，恐亦惊其子孙今日所考之伟绩。历代帝王经营天下，尚须甲兵钱粮；今之论史者开疆拓土，惟凭唇舌文字。其功之盛，岂不壮哉！

然世之奇谈，初起于笑谑，继而成于传闻，终乃奉若圭臬。传者不问其所自，闻者不究其所据。甲述于前，乙引于后，丙复引乙以证甲。辗转数番，遂成铁案。其间或有异议，则群起斥之，谓其不识祖宗，不知故土，不明大义。于是考史之学，渐成辨忠奸之术；求真之途，反为争信仰之场。

余观诸说，或荒诞绝伦，或诡谲莫测，或慷慨悲歌，或令人失笑。然其言既传于世，其人既深信不疑，则亦不妨录之。盖一代有一代之史，亦有一代之妄；正史存兴亡之迹，野史载闾巷之谈，而此等奇论，或足以见丙午世道人心之一斑。

故广搜屏间遗论，旁采坊间奇谈，汇而编之，名曰《丙午稗史》。不敢谓其为信史，亦不忍使其尽湮。凡所载者，悉存其说，以俟后人考辨。

他年若有好古之士，偶启此编，见其所论，必将疑曰：「此书所记，果为二十一世纪之人所言乎？」

余亦不知当何以答之。

惟愿后之览者，知此非上古神话，非志怪小说，非稗官杜撰，乃丙午之年，世间确有人如此言之、如此传之、如此信之。

\textbf{史或可伪，言终有迹；山河依旧，人心可书。}

是为序。


